\rtmtight
\begin{tblr}{width=\textwidth,colspec={|Q[c,m,2.1cm]|X[l,m]|},hlines,vlines,hspan=minimal,rowsep=1.5pt,colsep=3pt,row{1}={bg=headgray,font=\bfseries}}
\SetCell{c,m}\hfil\logo\hfil & \textbf{POLITEKNIK NEGERI MALANG}\par\textbf{JURUSAN TEKNOLOGI INFORMASI}\par\textbf{PROGRAM STUDI: D4 TEKNIK INFORMATIKA} \\
\SetCell[c=2]{l} \textbf{RENCANA TUGAS MAHASISWA (RTM) DAN RUBRIK PENILAIAN --- DRAF KONSEP C: SEIMBANG} & \\
Nama Mata Kuliah & Kuliah Kerja Nyata Tematik \\
Kode Mata Kuliah & RTI257002 (sementara) \quad \textbf{sks / Jam perminggu:} 2 SKS / 4 jam/minggu \quad \textbf{Semester:} 6/7 (sementara) \\
Dosen Pengampu & Tim Pengajar Program Studi D4 TI \\
\end{tblr}
\assessmentblock{Proposal Program Kerja KKN}{Case Method dan penyusunan dokumen}{SCPMK0102-06001, SCPMK0301-06002, SCPMK0803-06003}{Mahasiswa menyusun proposal program kerja KKN berdasarkan observasi dan analisis kebutuhan mitra, memuat rencana kegiatan pemberdayaan, jadwal, dan rancangan luaran KKN yang akan dihasilkan.}{Melakukan observasi lapangan, menganalisis kebutuhan mitra secara partisipatif, menyusun proposal, dan mempresentasikannya dalam seminar proposal.}{Dokumen proposal program kerja KKN dan bahan presentasi seminar proposal.}{Ketepatan analisis kebutuhan dan kelayakan program & 7 & Analisis kebutuhan mitra akurat dan mendalam; program kerja yang diusulkan realistis, relevan, dan menjawab masalah nyata mitra. & Analisis kebutuhan cukup memadai; program kerja relevan tetapi kurang detail pada beberapa aspek pelaksanaan. & Analisis kebutuhan dangkal atau tidak berbasis data lapangan; program kerja tidak relevan atau tidak realistis. \\ \\
Kejelasan rancangan luaran KKN & 7 & Rancangan luaran KKN (solusi TI dan/atau kegiatan) dijelaskan secara rinci, sesuai kebutuhan mitra, dan layak untuk diwujudkan pada tahap pelaksanaan. & Rancangan luaran cukup jelas tetapi beberapa bagian masih perlu disempurnakan. & Rancangan luaran tidak jelas, tidak sesuai kebutuhan mitra, atau tidak layak diwujudkan. \\ \\
Kualitas presentasi dan etika penyusunan & 6 & Presentasi proposal terstruktur dan meyakinkan; etika penyusunan (kejujuran data, pengakuan sumber) terjaga penuh. & Presentasi cukup baik dengan sedikit kekurangan struktur; etika penyusunan secara umum terjaga. & Presentasi tidak terstruktur atau terdapat pelanggaran etika penyusunan proposal. \\}

\assessmentblock{Pelaksanaan dan Keterlibatan}{Portofolio dan praktik lapangan}{SCPMK0102-06001, SCPMK0803-06003}{Mahasiswa melaksanakan program kerja KKN di lokasi mitra, mencatat aktivitas dalam logbook, dan menunjukkan keterlibatan aktif dalam kegiatan pemberdayaan maupun pengelolaan program.}{Mengisi logbook harian, melaksanakan kegiatan pemberdayaan dan pengelolaan program sesuai jadwal, serta mendapatkan validasi mitra/DPL secara berkala.}{Logbook digital/fisik yang terisi lengkap, divalidasi mitra/DPL, dan dokumentasi kegiatan pelaksanaan.}{Kelengkapan dan konsistensi pengisian logbook & 10 & Logbook terisi setiap kegiatan, mencakup aktivitas, waktu, hasil, dan refleksi; divalidasi mitra/DPL secara berkala. & Logbook terisi sebagian besar kegiatan dengan informasi yang cukup, tetapi beberapa entri tidak lengkap atau tidak tervalidasi. & Logbook banyak kegiatan kosong, isian tidak informatif, atau tidak divalidasi mitra/DPL. \\ \\
Tingkat keterlibatan dalam kegiatan pemberdayaan & 10 & Mahasiswa terlibat aktif dan konsisten dalam kegiatan pemberdayaan, memberikan kontribusi nyata yang diapresiasi mitra/komunitas. & Keterlibatan cukup aktif tetapi tidak konsisten sepanjang program; kontribusi ada namun terbatas. & Keterlibatan minim, mahasiswa pasif dalam kegiatan, atau tidak memberikan kontribusi nyata. \\ \\
Pengelolaan pelaksanaan dan penjadwalan program & 10 & Pelaksanaan program terjadwal dengan baik sesuai rencana kerja, kendala di lapangan dikelola secara proaktif. & Pelaksanaan cukup terjadwal; ada penyesuaian jadwal namun masih terkendali. & Pelaksanaan tidak terjadwal, banyak keterlambatan, atau tidak sesuai rencana kerja yang disepakati. \\}

\assessmentblock{Luaran KKN}{Portofolio dan praktik lapangan}{SCPMK0607-06004}{Mahasiswa merancang dan mewujudkan luaran KKN --- solusi TI dan/atau kegiatan pemberdayaan --- yang sesuai kebutuhan mitra dan diserahterimakan pada akhir program.}{Mendesain luaran, mengimplementasikan/melaksanakannya di lokasi mitra, mengujicobakan atau mendemonstrasikan kepada mitra, serta menyempurnakan berdasarkan umpan balik monev.}{Luaran KKN (solusi TI yang berjalan dan/atau laporan kegiatan pemberdayaan) beserta dokumentasi dan berita acara serah terima ke mitra.}{Kesesuaian luaran dengan kebutuhan mitra & 10 & Luaran KKN secara langsung menjawab kebutuhan mitra yang dirumuskan di proposal, disertai bukti pemanfaatan oleh mitra. & Luaran menjawab sebagian besar kebutuhan mitra; beberapa bagian belum optimal atau belum dimanfaatkan mitra. & Luaran tidak sesuai kebutuhan mitra yang telah disepakati atau tidak dapat dimanfaatkan. \\ \\
Kualitas rancang bangun/pelaksanaan luaran & 10 & Luaran dirancang dan diwujudkan dengan rapi dan terdokumentasi, menerapkan kaidah rekayasa atau pengelolaan kegiatan yang baik. & Luaran terwujud tetapi kualitas rancang bangun/pelaksanaan masih terbatas pada beberapa aspek. & Luaran tidak terwujud dengan baik atau tidak dapat dijelaskan secara teknis oleh mahasiswa. \\}

\assessmentblock{Laporan dan Seminar Hasil}{Case Method, penulisan ilmiah, dan presentasi}{SCPMK0301-06002, SCPMK0803-06003, SCPMK0607-06004}{Mahasiswa menyusun laporan KKN yang komprehensif serta mempresentasikan dan mempertahankan laporan beserta luaran KKN di hadapan tim penguji.}{Mengumpulkan data dan dokumentasi selama pelaksanaan program, menyusun laporan sesuai sistematika Polinema, bimbingan dengan DPL, revisi, dan presentasi seminar hasil.}{Naskah laporan KKN lengkap beserta lampiran (logbook, dokumentasi luaran, berita acara serah terima), dan bahan presentasi seminar hasil.}{Kelengkapan struktur dan dokumentasi pelaksanaan & 10 & Laporan memuat semua bab yang dipersyaratkan dengan dokumentasi pelaksanaan program dan luaran KKN yang lengkap dan akurat. & Struktur laporan lengkap tetapi beberapa bagian dokumentasi pelaksanaan/luaran kurang mendalam. & Struktur laporan tidak lengkap atau dokumentasi tidak mencerminkan kegiatan yang sesungguhnya. \\ \\
Kualitas presentasi dan komunikasi hasil & 10 & Presentasi terstruktur, disampaikan dengan percaya diri, dan mampu menyampaikan hasil program serta luaran KKN secara meyakinkan. & Presentasi cukup terstruktur dan dapat dipahami, tetapi penyampaian kurang percaya diri atau kurang meyakinkan. & Presentasi tidak terstruktur, bahasa tidak profesional, atau materi sulit dipahami penguji. \\ \\
Kemampuan menjawab pertanyaan dan keterkaitan luaran dengan kebutuhan mitra & 10 & Semua pertanyaan dijawab dengan tepat, mahasiswa dapat menjelaskan secara eksplisit keterkaitan luaran KKN dengan kebutuhan dan dampak nyata bagi mitra. & Sebagian besar pertanyaan dijawab cukup baik; keterkaitan luaran dengan kebutuhan mitra dijelaskan namun masih umum. & Banyak pertanyaan tidak dapat dijawab, atau tidak dapat menjelaskan keterkaitan luaran dengan kebutuhan mitra. \\}

\begin{tblr}{width=\textwidth,colspec={|Q[l,m,3.2cm]|X[l,m]|},hlines,vlines,hspan=minimal,row{1-3}={bg=headgray},rowsep=1.5pt,colsep=3pt}
Jadwal Pelaksanaan: & Minggu 1--16 sesuai RPS. Setiap evaluasi memiliki RTM dan rubrik multi-kriteria. \\
Lain-Lain Yang Diperlukan: & LMS, logbook digital, perangkat lapangan mitra, dan perangkat presentasi. \\
Pustaka: & Mengacu pustaka utama dan pendukung pada bagian RPS. \\
\end{tblr}
